% !TEX root = chapter-standalone.tex

\devel{
    \section{Motivation and Warmups}
    Since \SY{mappings between functors} may be rather abstract, we will first start with a concrete example.
    Imagine you are in Boston, and can go to New York or Providence.
    We may now define a category~$\CatC$ of cities, where morphisms are trips between the cities.
    Let~$\CatD$ be the category of pairs of real numbers under element-wise addition, and~$\funa, \funb\colon\CatC\fto\CatD$ be functors between them.
    $\funa$ gives us the distance and altitude difference between our starting point (Boston) and our desired destination, and~$\funb$ gives us the gas price relative to our starting point.

    Now, we can create a mapping~$\ntrafoa$ mapping the distances to gas prices.
    Thinking about it, it relates \SY{the elements of the category of distances to those of the category of prices}.
    If we write out our categories in typical diagrams, we notice that they \SY{commute} (see \cref{def:commutative-diagram}).

    \middlesag{diag_boston_nyc}

    In other words, for each pair of (distance, altitude) relative to Boston, we can compute the gas price of the trip.
}

\section{Natural transformations}
\linkvideo{spring2021-nat-trafos:diagrams} % Diagrams
%\linkvideo{spring2021-nat-trafos:natural-trafos} % Natural transformations

To formally define \SY{natural transformations}, the general situation we will start from is when we have two \SY{functors}~$\funa \colon \CatC \fto \CatD$ and~$\funb \colon \CatC \fto \CatD$, sharing the same source and target, respectively.

A \SY{natural transformation} $\ntrafoa$ from~$\funa$ to~$\funb$ is then a kind of ``translation'' from~$\funa$ to~$\funb$. How might one define such a thing?

First, let's look at the situation only on the level of objects.
Each object~$\Obja$ of~\CatC is mapped by~$\funa$ and~$\funb$ to an object~$\funa(\Obja)$ and~$\funb(\Obja)$ of~\CatD, respectively.
One straightforward way to relate~$\funa(\Obja)$ to~$\funb(\Obja)$ is to choose a morphism~$\funa(\Obja) \mto \funb(\Obja)$ in~\CatD.
We call this morphism~$\ntrafoa_{\Obja}$, using the subscript~$\Obja$ since~$\ntrafoa_{\Obja}$ relates the respective images of the object~$\Obja$ under~$\funa$ and~$\funb$.

\begin{figure*}[h]
    \centering
    \begin{ctdefinitionshade}
        \small
        \includesag{094_natural_setup_obj}
    \end{ctdefinitionshade}
    \caption{}
    \label{fig:natural_setup_obj}
\end{figure*}

If we choose such a morphism for each object in~\CatC, then we have collection~$\makeset{ \ntrafoa_\Obja }_{\Obja \setin \ObC}$ of morphisms in~\CatD, indexed by the objects of~\CatC.

Next, consider a morphism~$\mora \colon \Obja \mto \Objb$ in the category~\CatC.
Under the \SY{functor}~$\funa$ it will be mapped to some morphism~$\funa(\mora)\colon \funa(\Obja) \mto \funa(\Objb)$ in~\CatD, and under the \SY{functor}~$\funb$ it will be mapped to some other morphism~$\funb(\mora)\colon \funb(\Obja) \mto \funb(\Objb)$, also in~\CatD. We can think of~$\funa(\mora)\colon \funa(\Obja) \mto \funa(\Objb)$ and~$\funb(\mora)\colon \funb(\Obja) \mto \funb(\Objb)$ as each being (very small) diagrams in~\CatD.
\begin{figure*}[h!]
    \centering
    \begin{ctdefinitionshade}
        \small
        \includesag{095_natural_setup_mor}
    \end{ctdefinitionshade}
    \caption{}
    \label{fig:natural_setup_mor}
\end{figure*}
%\begin{marginfigure}
%    \hfill
%    \includesag{diagram_two_functors}
%    % \caption{}
%    % \label{fig:diagram_two_functors}
%\end{marginfigure}




%\begin{marginfigure}
%    \hfill
%    \includesag{diagram_two_functors_conn}
%    % \caption{}
%    % \label{fig:two-functors-connected-diagram}
%\end{marginfigure}


If we have already chosen morphisms~$\ntrafoa_\Obja\colon \funa(\Obja) \mto \funb(\Obja)$ and~$\ntrafoa_\Objb\colon \funa(\Objb) \mto \funb(\Objb)$ in~\CatD, then these will connect the two diagrams in $\CatD$.

\begin{figure*}[h!]
    \centering
    \begin{ctdefinitionshade}
        \small
        \includesag{096_natural_setup_square}
    \end{ctdefinitionshade}
    \caption{}
    \label{fig:natural_setup_square}
\end{figure*}



% \equationsag{diagram_two_functors_conn}{eq:two-functors-connected-diagram}
% \begin{figure}[h!]
%     \centering
%     \includesag{diagram_two_functors_conn}
%     \caption{}
%     \label{fig:two-functors-connected-diagram}
% \end{figure}

%
\begin{marginfigure}
    \hfill
    \includesag{naturality_square}
    \caption{}
    \label{fig:naturality-square}
\end{marginfigure}

In~\CatD, this gives rise to the square diagram shown to the side.
We'll require, as a condition on the morphisms~$\ntrafoa_\Obja$ and~$\ntrafoa_\Objb$, that they make the diagram commutative:
\begin{equation}
    \funa(\mora) \mthen \ntrafoa_\Objb = \ntrafoa_\Obja \mthen \funb(\mora).
\end{equation}

Now consider not only a single morphism~$\mora \colon \Obja \mto \Objb$ in~\CatC being mapped by~$\funa$ and~$\funb$, respectively, but \emph{all} of the category~\CatC.
Under~$\funa$, the category~\CatC is mapped to a -- possibly very complicated -- diagram in~\CatD (a directed graph of objects and morphism comprising the image of~$\funa$), and similarly, under~$\funb$, the category~\CatC is mapped to another diagram in~\CatD (the image of~$\funb$).

To relate the image of~$\funa$ to the image of~$\funb$ we can proceed in the same way as above: for each object~$\Obja$ in~\CatC, we choose a morphism~$\ntrafoa_\Obja \colon \funa(\Obja) \mto \funb(\Obja)$ in~\CatD.
In other words, we have a collection of morphisms~$(\ntrafoa_\Obja)$, ${\Obja \setin \ObC}$ indexed by the objects of~\CatC.
These gives rise to lots of squares of the kind in \cref{fig:naturality-square}, which we will require to be commutative.
It is because of this commutativity condition, which is a condition on the collection~$(\ntrafoa_\Obja)$, ${\Obja \setin \ObC}$, that some mathematicians would say that the collection~$(\ntrafoa_\Obja)$, ${\Obja \setin \ObC}$ is a ``coherent'' or ``natural'' way to relate the image of~$\funa$ to the image of~$\funb$.
(This does not mean, however, that there is at most one \SY{natural transformation} between any two given \SY{functors}---on the contrary, there might be many!)

\begin{marginfigure}
    \centering
    \includesag{two_naturality_squares}
    \caption{}
    \label{fig:two-naturality-squares}
\end{marginfigure}

In \cref{fig:two-naturality-squares} we have illustrated a situation involving three objects and two morphisms in~\CatC, giving rise to two squares.
We have ``glued'' the two squares together since they share an edge (this a more compact way of drawing them).
%Note that because each of the two component squares in the diagram commute, so does the entire diagram.

\linkvideo{spring2021-nat-trafos:natural-trafos:nat-trafo-def} % Definition of natural transformation

\begin{ctdefinition}[Natural transformation]
    \label{def:natural-transformation}
    Let \CatC and \CatD be categories, and let~$\funa,\funb\colon \CatC\fto \CatD$ be \SY{functors}.
    A \maindef{natural transformation}~$\ntrafoa \colon \funa \nto \funb$ is specified by:

    \constit
    \begin{enumerate}
        \item For each object~$\Obja\setin \ObC$, a morphism $\ntrafoa_\Obja \colon \funa(\Obja)\mto \funb(\Obja)$ in \CatD, called the $\Obja$\emph{-component} of $\ntrafoa$.
    \end{enumerate}
    \condit
    \begin{enumerate}
        \item For every morphism~$\mora\colon \Obja\mto \Objb$ in \CatC, the components of $\ntrafoa$ must satisfy the \emph{naturality condition}
              \begin{equation}
                  \funa(\mora)\mthen \ntrafoa_\Objb = \ntrafoa_\Obja\mthen \funb(\mora).
              \end{equation}
              In other words, the following diagram must commute:
              \begin{equation}
                  \label{eq:def-naturality-square}
                  \middlesag{55_natural_2}
              \end{equation}
    \end{enumerate}
\end{ctdefinition}

To reiterate: a \SY{natural transformation} $\ntrafoa$ is a \emph{collection} $(\ntrafoa_\Obja)_{\Obja \setin \ObC}$ of morphisms (called the \emph{components} of the natural transformation) which satisfy the naturality conditions.
The name ``components'' is analogous to how a vector $v = (v_1, .
    .., v_n)$ has \emph{components}, or to how a sequence $a = (a_n)_{n \setin \natnumbers}$ has \emph{terms}.

The diagrams \cref{eq:def-naturality-square} are often called \emph{naturality squares}, and a \SY{natural transformation}~$\ntrafoa \colon \funa\nto \funb$ is often depicted concisely in this manner:
\equationsag{55_natural_1}{eq:55_natural_1}

\Cref{fig:nat_trans_graphically} shows a diagram that describes the property of functors and of natural transformations.
The diagram in $\CatD$ is a ``commuting prism'': all faces of the prism commute.

\begin{figure*}[h]
    \centering
    \begin{ctdefinitionshade}
        \small
        \includesag{096_natural_graphically}
    \end{ctdefinitionshade}
    \caption{}
    \label{fig:nat_trans_graphically}
\end{figure*}

%But perhaps this prism is better explained with the help of some examples:

\subsection{Some initial examples from mathematics}

\begin{example}
    Consider the following two functors $\funa, \funb \colon \Set \Ctimes \Set \Ctimes \Set \fto \Set$.
    We define $\funa$ on objects by
    \begin{equation}
        \funa(\tup{\setA, \setB, \setC}) = (\setA \cartprod \setB) \cartprod \setC
    \end{equation}
    and define $\funb$ on objects by
    \begin{equation}
        \funb(\tup{\setA, \setB, \setC}) = \setA \cartprod (\setB \cartprod \setC).
    \end{equation}
    For their actions on morphisms, consider a morphism
    \begin{equation}
        \tup{\mora, \morb, \morc} \colon \tup{\setA, \setB, \setC} \mto \tup{\setAprime, \setBprime, \setCprime}
    \end{equation}
    in $\Set \Ctimes \Set \Ctimes \Set$.
    Its image under $\funa$ is
    \begin{equation}
        \tup{\tup{\mora, \morb}, \morc} \colon (\setA \cartprod \setB) \cartprod \setC \mto (\setAprime \cartprod \setBprime) \cartprod \setCprime
    \end{equation}
    and its image under $\funb$ is
    \begin{equation}
        \tup{\mora, \tup{\morb, \morc}} \colon \setA \cartprod (\setB \cartprod \setC) \mto \setAprime \cartprod (\setBprime \cartprod \setCprime).
    \end{equation}

    One way to see that~$\funa$ (and similarly $\funb$) is indeed a functor is to note that it is equal to the following composition of functors
    \begin{equation}
        \begin{aligned}
            \Set \Ctimes \Set \Ctimes \Set & \fto \Set \Ctimes \Set \fto \Set, \\
            \tup{\setA, \setB, \setC}      & \mapsto \tup{\setA \cartprod \setB, \setC} \mapsto (\setA \cartprod \setB) \cartprod \setC
        \end{aligned}
    \end{equation}
    and recall from \cref{ex:CartProdAsFunctor} that ``$\cartprod$'' is a functor.

    Now we define a natural transformation $\ntrafoa \colon \funa \nto \funb$ by specifying its components to be the functions
    \begin{equation}
        \defmapcomma{ \ntrafoa_{\tup{\setA, \setB, \setC}}}{(\setA \cartprod \setB) \cartprod \setC}{\fto}{\setA \cartprod (\setB \cartprod \setC)}{\tup{\tup{\ela, \elb}, \elc}}{\tup{\ela, \tup{\elb, \elc}}}
    \end{equation}
    indexed by triples of sets $\tup{\setA, \setB, \setC}$.

    For the family of morphisms $\ntrafoa_{\tup{\setA, \setB, \setC}}$ to be a natural transformation, we need to check that the diagrams
    %
    \equationsag{ntrafo_associator_cart_prod}{eq:ntrafo_associator_cart_prod}
    %
    in $\Set$ commute for all morphisms $\tup{\mora, \morb, \morc}$ in $\Set \Ctimes \Set \Ctimes \Set$.
    It is easily checked that this is true.

    This natural transformation is an example of something called an \emph{associator}, which we will discuss later when we define monoidal categories.
    The idea here is that the cartesian product of sets is not quite an associative operation, but almost: instead of an ``equality'' symbol in the usual equation for the associative law, we have the components of this associator natural transformation.
\end{example}

\todotext{J: the last paragraph of the above example needs to be edited to reflect that we already introduced associators in the chapter on sets and functions}

\begin{example}
    Let $\funa \colon \Set \Ctimes \Set \fto \Set$ be the functor which on objects maps any pair of sets $\tup{\setA, \setB}$ to their cartesian product $\setA \cartprod \setB$.
    On functions, it maps any pair of functions $\tup{\mora, \morb}$ to their cartesian product $\mora \funcprod \morb$.

    Consider another functor $\funb \colon \Set \Ctimes \Set \fto \Set$ which on objects maps any pair of sets $\tup{\setA, \setB}$ to their cartesian product $\setB \times \setA$ (and similarly for morphisms).
    In other words, $\funb$ is very similar to $\funa$, however it is different in that, when forming the cartesian product, the order of the factors is swapped.

    There is a natural transformation $\ntrafoa \colon \funa \nto \funb$ which expresses explicitly the relationship between $\funa$ and $\funb$.
    Its components are these functions:
    \begin{equation}\label{eq:symmetry-nat-trafo}
        \defmapperiod{\ntrafoa_{\tup{\setA, \setB}}}{\setA \cartprod \setB}{\mto }{\setB \cartprod \setA}{\tup{\ela, \elb}}{\tup{\elb, \ela}}
    \end{equation}

    %Let us now check that this family of functions $\ntrafoa_{\tup{\setA, \setB}}$ does indeed define a natural transformation. For this we need to check that the naturality squares all commute. To check this, suppose $\tup{\mora, \morb} \colon \tup{\setA, \setB} \mto \tup{\setC, \setD}$ is an arbitrary morphism in $\Set \Ctimes \Set$.
\end{example}

\begin{gradedexercise}[\exname{NaturalBraiding}]

    Consider the functor $\funa \colon \Set \Ctimes \Set \fto \Set$ defined on objects and morphisms by
    \begin{equation}
        \tup{\setA, \setB} \mapsto \setA \cartprod \setB \qquad \qquad \tup{\mapa, \mapb} \mapsto \mapa \funcprod \mapb,
    \end{equation}
    and consider as well the similar (but different!) functor  $\funb \colon \Set \Ctimes \Set \fto \Set$ defined on objects and morphisms by
    \begin{equation}
        \tup{\setA, \setB} \mapsto \setB \cartprod \setA \qquad \qquad \tup{\mapa, \mapb} \mapsto \mapb \funcprod \mapa.
    \end{equation}

    Now consider, for each ordered pair of sets $\tup{\setA, \setB}$, the function
    \begin{equation}\label{eq:braiding-components}
        \defmapperiod{
            \braiding_{\tup{\setA, \setB}}
        }{
            \setA \cartprod \setB
        }{
            \to
        }{
            \setB \cartprod \setA
        }{
            \tup{\ela, \elb}
        }{
            \tup{\elb, \ela}
        }
    \end{equation}

    Your task is to prove that the family of functions $\{ \braiding_{\tup{\setA, \setB}} \}_{\tup{\setA, \setB}}$ defines a natural transformation $\braiding \colon \funa \nto \funb$.

\end{gradedexercise}

\solutionof{NaturalBraiding}

\begin{example}
    Consider the \SY{powerset} \SY{functor} $\powerset$ from \cref{ex:powerset_functor}.
    As a reminder, the \SY{powerset} \SY{functor} $\powerset$ maps a set $\setA$ to its \SY{powerset} $\powerset(\setA)$, and a function $\mora \colon \setA \mto \setB$ to the function $\powerset(\mora) \colon \powerset(\setA) \mto \powerset(\setB)$ which sends each subset of $\setA$ to its image under $\mora$, which is a subset of $\setB$.

    There is a \SY{natural transformation}~$\ntrafoa\colon \funid_\Set \nto \powerset$ whose components are the functions
    \begin{equation}
        \begin{aligned}
            \ntrafoa_\setA\colon \setA & \mto \powerset(\setA) \\
            \setAel                    & \mapsto \makeset{\setAel}.
        \end{aligned}
    \end{equation}
    In other words, the \SY{natural transformation} embeds each element of~\setA into the power set~$\powerset(\setA)$.
    To check that this is a \SY{natural transformation}, consider an arbitrary function between sets $\mora\colon \setA\mto \setB$.
    On the one hand,
    \begin{equation}
        (\ntrafoa_\setA \mthen \powerset(\mora))(\setAel)
        = \powerset(\mora)(\makeset{\setAel}) =\makeset{\mora(\setAel)},
    \end{equation}
    while on the other hand
    \begin{equation}
        (\mora \mthen \ntrafoa_\setB)(\setAel) = \ntrafoa_\setB (\mora(\setAel)) = \makeset{\mora(\setAel)}.
    \end{equation}
    Thus, the condition for $\ntrafoa$ to be a natural transformation is satisfied.
\end{example}

\todotext{J: Make a diagram for the above example}

\begin{gradedexercise}[\exname{ListUnitNatural}]

    Recall the list functor $\listfun \colon \Set \fto \Set$ which, on objects, assigns to each set $\setA$ the set $\listfun(\setA)$ of finite lists in elements of $\setA$, and to each function $\mapa \colon \setA \to \setB$ it assigns the induced function
    \begin{equation}\label{eq:list-functor-on-functions}
        \defmapperiod{
            \listfun(\mapa)
        }{
            \listsof{\setA}
        }{
            \to
        }{
            \listsof{\setB}
        }{
            \makegenlist{ \setAel_i }
        }{
            \makegenlist{ \mapa(\setAeln{i}) }
        }
    \end{equation}
    (You do not need to prove that this is a functor; we take that fact as given in this exercise.)

    Consider now the family of functions $\{ \ntrafoa_\setA \}_\setA$ defined by
    \begin{equation}
        \defmapperiod{
            \ntrafoa_\setA
        }{
            \setA
        }{
            \to
        }{
            \listsof{\setA}
        }{
            \ela
        }{
            \maketypedlist{\ela}{\setA}
        }
    \end{equation}
    In other words, the function $\ntrafoa_\setA$ maps each element $\ela$ of $\setA$ to a list of length $1$ whose single entry is the element $\ela$.
    Does this family of functions define a natural transformation? Prove your answer.
\end{gradedexercise}

\solutionof{ListUnitNatural}

\clearpage

\section{More examples}
\linkvideo{spring2021-nat-trafos:natural-trafos:double-dual} % Double dual

\todojira{631}{\alphubel: @JL: write up example here involving monoid actions / state machines}

\todojira{632}{\alphubel: @JL: write up example here about natural transformations between monotone maps}


\begin{example}
    Fix a set $\subA$.
    There are functors $\funa, \funb: \Set\op \Ctimes \Set \fto \Set$ whose respective actions on objects are
    \begin{equation}
        \funaob \colon \tup{\setA, \setB} \mapsto \Hom_\Set(\setA \cartprod \subA, \setB)
    \end{equation}
    and
    \begin{equation}
        \funbob \colon \tup{\setA, \setB} \mapsto \Hom_\Set(\setA, \setB^\subA).
    \end{equation}
    These functors may be understood as built up using compositions of functors of the kind discussed in \cref{ex:MultiplicationWithASet}, \cref{ex:ExponentiationWithASet} and \cref{exa:hom-functor}.

    Recall that we can ``curry'' any function $\mora : \setA \cartprod \subA \mto \setB$ to get a function $\hat \mora \colon \setA \mto \setB^\subA$, where $\hat \mora(\ela)$ may be thought of as a partial evaluation of $\mora$.

    There is a natural transformation $\ntrafoa \colon \funa \nto \funb$ whose components are the functions
    \begin{equation}\label{eq:currying-nat-trafo}
        \ntrafoa_{\tup{\setA, \setB}} \colon \Hom_\Set(\setA \cartprod \subA, \setB) \mto \Hom_\Set(\setA, \setB^\subA), \ \mora \mapsto \hat \mora,
    \end{equation}
    where $\hat \mora$ is the ``curried'' version of $\mora$.
\end{example}

\vfill

\begin{marginfigure}
    \centering
    \includesag{graph-cat}
    \caption{}
    \label{fig:graph-cat-again}
\end{marginfigure}

\begin{gradedexercise}[\exname{NatTrafosGraphs}]
    \label{ex:NatTrafosGraphs}
    This exercise builds on \cref{ex:GraphsViaFunctors}.
    There, we defined a category~$\CatSymbol{G}$ which has precisely two objects and four morphisms, see \cref{fig:graph-cat-again} (the two \SY{identity morphisms} are not drawn).
    The task there was to understand how specifying a \SY{functor} from this category $CatSymbol{G}$ into the category of sets is ``the same thing'' as specifying a directed graph.

    Now consider two \SY{functors}~$\funaA, \funaB \colon \CatSymbol{G} \fto \Set$.
    Spell out what it means to have a \SY{natural transformation}~$\ntrafoa\colon \funaA \nto \funaB$.
    What does this correspond to in the language of directed graphs?
\end{gradedexercise}

\solutionof{NatTrafosGraphs}

\begin{gradedexercise}[\exname{UpperSetsNatTrafos}]
    \label{ex:UpperSetsNatTrafos}
    This exercise builds on \cref{ex:UpperSetsViaFunctors}.
    There we fixed a poset~\posA, viewed it as a category~$\CatSymbol{P}$, and saw that \SY{functors}~$\CatSymbol{P} \fto \Bool$ encode \SY{upper sets} in~\posA.
    Suppose we have two \SY{functors}~$\funaA, \funaB \colon \CatSymbol{P} \fto \Bool$.
    What does a \SY{natural transformation}~$\ntrafoa \colon \funaA \nto \funaB$ correspond to in terms of the \SY{upper sets} encoded by~$\funaA$ and~$\funaB$, respectively?
\end{gradedexercise}

\solutionof{UpperSetsNatTrafos}



\begin{gradedexercise}[\exname{NaturalCurry}]
    \label{ex:NaturalCurry}
    Fix a set $\subA$.
    There are functors $\funa, \funb: \Set\op \Ctimes \Set \fto \Set$ whose respective actions on objects are
    \begin{equation}
        \funaob \colon \tup{\setA, \setB} \mapsto \Hom_\Set(\setA \cartprod \subA, \setB)
    \end{equation}
    and
    \begin{equation}
        \funbob \colon \tup{\setA, \setB} \mapsto \Hom_\Set(\setA, \setB^\subA).
    \end{equation}

    On morphisms, $\funa$ acts as follows.
    Given a morphism $\tup{\mora\op, \morb} \colon \tup{\setA, \setB} \mto \tup{\setAprime, \setBprime}$ in $\Set\op \Ctimes \Set$, the function
    $$\funa(\tup{\mora\op, \morb}) \colon \Hom_\Set(\setA \cartprod \subA, \setB) \mto \Hom_\Set(\setAprime \cartprod \subA, \setBprime)$$
    takes any function $\stylemorph{\varphi} \colon \setA \cartprod \subA \mto \setB$ and maps it to the function
    \begin{equation}
        \tup{\mora, \catid_{\subA}} \mthen \stylemorph{\varphi} \mthen \morb \colon \setAprime \cartprod \subA \mto \setBprime.
    \end{equation}

    Let us also define the action of $\funb$ on morphisms.
    Given again a morphism $\tup{\mora\op, \morb} \colon \tup{\setA, \setB} \mto \tup{\setAprime, \setBprime}$ in $\Set\op \Ctimes \Set$, the function
    $$\funb(\tup{\mora\op, \morb}) \colon \Hom_\Set(\setA, \setB^\subA) \mto \Hom_\Set(\setAprime, \setBprime^\subA)$$
    takes any function $\stylemorph{\varphi} \colon \setA \mto \setB^\subA$ and maps it to the function
    \begin{equation}
        \mora \mthen \stylemorph{\varphi} \mthen \morb_{\stylemorph{*}} \colon \setAprime \mto \setBprime^\subA
    \end{equation}
    where $\morb_{\stylemorph{*}}$ is the function
    \begin{equation}
        \morb_{\stylemorph{*}} \colon \setB^\subA \mto \setBprime^\subA, \ \stylemorph{\psi} \mapsto \stylemorph{\psi} \mthen \morb.
    \end{equation}

    %These functors may be understood as built up using compositions of functors of the kind discussed in \cref{ex:MultiplicationWithASet}, \cref{ex:ExponentiationWithASet} and \cref{exa:hom-functor}.

    Now, recall that we can ``curry'' any function $\mora : \setA \cartprod \subA \mto \setB$ to get a function $\stylemorph{\overline{{\mora}}} \colon \setA \mto \setB^\subA$, where $\stylemorph{\overline{{\mora}}}(\ela)$ may be thought of as a partial evaluation of $\mora$.

    Your task in this exercise: show that the functions
    \begin{equation}
        \ntrafoa_{\tup{\setA, \setB}} \colon \Hom_\Set(\setA \cartprod \subA, \setB) \mto \Hom_\Set(\setA, \setB^\subA), \ \mora \mapsto \stylemorph{\overline{{\mora}}},
    \end{equation}
    where $\stylemorph{\overline{{\mora}}}$ is the ``curried'' version of $\mora$, are the components of a natural transformation $\ntrafoa \colon \funa \nto \funb$.
\end{gradedexercise}

\solutionof{NaturalCurry}

\devel{

\begin{example}
    \label{ex:Vect}
    Consider the category~$\VectR$ whose objects are real \SY{vector spaces} and whose morphisms are linear maps.
    (For convenience, in the following we sometimes omit reference to the ground field.)
    Recall that the \emph{dual} of a \SY{vector space}~$V$ is the \SY{vector space} describing all linear maps from~$V$ to~\reals:
    \begin{equation}
        \label{eq:dual-vector-space}
        V^* \definedas \HomSet{\VectR}{V}{\reals},
    \end{equation}
    Also, recall that if~$\mapa\colon V \mto W$ is a linear map, then its dual is the linear map~$\mapa^* \linebreak[1] \colon W^* \mto V^*$ which maps any $\xi \setin W^*$ to the element of $V^*$ given by
    \begin{equation}
        \mapa^*(\xi) \colon V \mto \reals, v \mapsto \xi(\mapa(v)).
    \end{equation}

    Applying the above duality construction twice to a \SY{vector space} or a linear map gives their double dual.
    It turns out that this is a functorial operation.
    That is, there is a functor
    \begin{equation}
        \label{eq:double-dual-functor}
        \text{\stylefunctors{Double dual}}\colon \VectRs \fto \VectRs
    \end{equation}
    that maps every \SY{vector space} and every linear map to its double dual.

    Furthermore, for any \SY{vector space}~$V$, there is a ``canonical'' or ``natural'' map
    \begin{equation}
        \ntrafoa_V \colon V \nto V^{**}
    \end{equation}
    defined by
    \begin{equation}
        \label{eq:natural-trafo-to-double-dual}
        \ntrafoa_V(v)(l) = l(v), \quad v \setin V, l \setin V^*.
    \end{equation}
    These form the components of a \SY{natural transformation} from the \SY{identity functor} on $\VectRs$ to the double dual \SY{functor}.
    \equationsag{nat-trafo-ddual}{eq:nat-trafo-ddual}
\end{example}

\begin{gradedexercise}[\exname{DoubleDualNatTrafo}]
    \label{ex:DoubleDualNatTrafo}
    This exercise builds on \cref{ex:DoubleDualFunctor}.
    Consider the category~$\VectR$ whose objects are real \SY{vector spaces} and whose morphisms are linear maps.

    For any \SY{vector space}~$V$, there is a ``canonical'' or ``natural'' map
    \begin{equation}
        \ntrafoa_V \colon V \nto V^{**}
    \end{equation}
    defined by
    \begin{equation}
        \ntrafoa_V(v)(l) = l(v), \quad v \setin V, l \setin V^*.
    \end{equation}

    Your tasks in this exercise:
    \begin{enumerate}
        \item Check that the operation of ``taking the double dual'' (see \cref{eq:double-dual-functor} above) defines a functor.
        \item Verify that these components define a \SY{natural transformation} from the \SY{identity functor} on $\VectRs$ to the double dual \SY{functor}.
    \end{enumerate}
    \equationsag{nat-trafo-ddual}{eq:ex:nat-trafo-ddual}
\end{gradedexercise}

\solutionof{DoubleDualNatTrafo}

}


\devel{

    %\section{Horizontal composition}

    %\todojira{638}{\bernina: @JL: write this up}

    \equationsag{horizontal_composition_1}{eq:horizontal_composition_1}

    \equationsag{horizontal_composition_2}{eq:horizontal_composition_2}

    \section{To add}
    \todojira{141}{\bernina: @JL: Add parts corresponding to listed videos.
    }
    %\linkvideo{spring2021-nat-trafos:natural-trafos:horizontal-composition} % Horizontal composition
    %\linkvideo{spring2021-nat-trafos:natural-trafos:interchange-law} % Interchange law
    \linkvideo{spring2021-nat-trafos:natural-trafos:rel-mon-maps} % Relating monotone maps
    \linkvideo{spring2021-nat-trafos:natural-trafos:eq-maps-gr-actions} % Equivariant maps between group actions
    
    \todotext{Add section or subsection somewhere in this chapter that explains how natural transformations can be used to encode the 'cones' and 'co-cones' that appear in the definition of limit and colimit}
    
}




\section{Composition of natural transformations}
\linkvideo{spring2021-nat-trafos:natural-trafos:nattrafos-as-mor} % Natural transformations are morphisms between functors
We have seen that natural transformations between two functors $\funa, \funb \colon \CatC \fto \CatD$ map objects in $\CatC$ to morphisms in $\CatD$, and map morphisms in $\CatC$ to commutative diagrams.
This is quite similar to the effect of functors on category objects and morphisms.


What if there were a category where objects are functors from $\CatC$ to $\CatD$, and morphisms are natural transformations?

\subsection{Vertical Composition}
As for any category, we would need to define the morphism composition law and the identity morphisms.
The former would look like this:

\equationsag{naturality_vert_comp}{eq:naturality-vert-comp}

Due to its diagrammatic form, we call this type of composition \SY{vertical composition}.
\begin{ctdefinition}[Vertical Composition of natural transformations]\label{def:composition-of-naturali-tranformations}
    Let $\CatC, \CatD$ be categories and let $\funa, \funb, \func \colon \CatC \fto \CatD$ be functors from $\CatC$ to $\CatD$.
    Suppose we are given natural transformations
    \begin{equation}
        \ntrafoa \colon \funa \nto \funb,
    \end{equation}
    \begin{equation}
        \ntrafob \colon \funb \nto \func.
    \end{equation}
    Their (vertical) composition $\ntrafoa \nthen \ntrafob$ is the natural transformation
    \begin{equation}
        \ntrafoa \nthen \ntrafob \colon \funa \nto \func
    \end{equation}
    defined in components by
    \begin{equation}\label{eq:nat-trafo-composition}
        (\ntrafoa \nthen \ntrafob)_\Obja \definedas \ntrafoa_\Obja \mthen \ntrafob_\Obja \quad \quad \forall \ \Obja \setin \Ob_\CatC.
    \end{equation}
\end{ctdefinition}

\begin{ctdefinition}[Identity natural transformation]\label{def:identity-natural-transformation}
    Let $\CatC, \CatD$ be categories and let $\funa \colon \CatC \fto \CatD$ be a functor.
    The identity natural transformation at $\funa$ is the natural transformation $\natid_\funa \colon \funa \nto \funa$ defined in components by
    \begin{equation}\label{eq:identity-nat-trafo-def}
        (\natid_\funa)_\Obja \definedas \catid_{\funa(\Obja)} \quad \quad \forall \ \Obja \setin\Ob_\CatC.
    \end{equation}
\end{ctdefinition}

\begin{ctdefinition}\label{def:category-of-functors}
    Let $\CatC, \CatD$ be categories.
    The category $[\CatC, \CatD]$ of functors from $\CatC$ to $\CatD$ is given by
    \begin{enumerate}
        \item \emph{Objects}: functors $\CatC \fto \CatD$.
        \item \emph{Morphisms}: natural transformations between functors $\CatC \fto \CatD$.
        \item \emph{Composition}: (vertical) composition of natural transformations.
        \item \emph{Identities}: identity natural transformations.
    \end{enumerate}
\end{ctdefinition}

\subsection{Natural isomorphisms}

\begin{ctdefinition}[Natural isomorphism]
    \label{def:natural-isomorphism}
    A \SY{natural transformation}~$\ntrafoa \colon \funa \nto \funb $ is called a \maindef{natural isomorphism} if each component morphism ~$\ntrafoa_\Obja$ in \CatD is an isomorphism.
\end{ctdefinition}

\begin{lemma}
    Let $\CatC, \CatD$ be categories and let $\funa, \funb \colon \CatC \fto \CatD$ be functors.
    A natural transformation $\ntrafoa \colon \funa \nto \funb $ is an isomorphism in the category of functors $[\CatC, \CatD]$ if and only if $\ntrafoa$ is a natural isomorphism.
\end{lemma}

\subsection{Horizontal Composition}
\linkvideo{spring2021-nat-trafos:natural-trafos:horizontal-composition} % Horizontal composition

Suppose we have three categories~$\CatC, \CatD$ and~$\CatE$, functors~$\funaA,\funbA \colon\CatC \fto \CatD$ and~$\funaB,\funbB\colon \CatD \fto \CatE$, and also two natural transformations
$\ntrafoa_\stylenat{1} \colon \funaA \nto \funbA$ and~$\ntrafoa_\stylenat{2} \colon \funaB \nto \funbB$.
We would then have the following situation.

\equationsag{naturality_horizontal}{eq:naturality-horizontal}

This looks suspiciously composable.

\begin{ctdefinition}[Horizontal Composition of natural transformations]
    \label{def:horizontal-composition}
    Let~$\CatC, \CatD, \CatE$ be categories and let~$\funaA,\funbA \colon\CatC \fto \CatD$ and $\funaB,\funbB\colon \CatD \fto \CatE$ be functors.
    Suppose we are given natural transformations
    \begin{equation}
        \ntrafoa_\stylenat{1}  \colon \funaA \nto \funbA, \quad \text{ and } \quad     \ntrafoa_\stylenat{2}  \colon \funaB \nto \funbB.
    \end{equation}
    Their horizontal composition is the natural transformation
    \begin{equation}
        \ntrafoa_\stylenat{1}  \nthenhor \ntrafoa_\stylenat{2}  \colon (\funaA \fthen \funbA) \nto (\funaB \fthen \funbB)
    \end{equation}
    defined in components by
\begin{equation}
({\ntrafoa_\stylenat{1} \stylenat{*} \ntrafoa_\stylenat{2}})_{\Obja} = {\ntrafoa_\stylenat{2}}_{\funa_\stylefunctors{1}(\Obja)} \mthen \funb_\stylefunctors{2}({\ntrafoa_\stylenat{1}}_{\Obja}).
\end{equation}

    In other words:
    \begin{equation}
            ({\ntrafoa_\stylenat{1} \stylenat{*} \ntrafoa_\stylenat{2}})_{\Obja} \colon (\funa_\stylefunctors{2}
            ( \funa_\stylefunctors{1}(\Obja)) \overset{{\ntrafoa_\stylenat{2}}_{\funa_\stylefunctors{1}(\Obja)}}{\mto} (\funb_\stylefunctors{2}
            ( \funa_\stylefunctors{1}(\Obja)) \overset{\funb_\stylefunctors{2}({\ntrafoa_\stylenat{1}}_{\Obja})}{\mto} (\funb_\stylefunctors{2} (\funb_\stylefunctors{1}(\Obja)).
    \end{equation}
\end{ctdefinition}

\todotext{J: make remark here that horizontal composition could also have been defined in a different way but equivalent way, thanks to the interchange law... }

\subsection{Interchange Law}
\linkvideo{spring2021-nat-trafos:natural-trafos:interchange-law} % Interchange law

The following statement is quite powerful.
Although perhaps obvious-looking, it allows us to presume associativity of
natural transformations.

\begin{proposition}[Interchange] \label{prop:interchange}
    Let~$\CatC, \CatD, \CatE$ be categories, $\funaA, \funbA, \funcA\colon\CatC \fto \CatD$, $\funaB, \funbB, \funcB:\CatD \fto \CatE$ be functors and $\ntrafoaA\colon\funaA\nto\funbA$, $\ntrafoaB\colon\funaB\nto\funbB$, $\ntrafobA\colon\funbA\nto\funcA$, $\ntrafobB\colon\funbB\nto\funcB$ be natural transformations.
    Then,
    \begin{equation}
        (\ntrafoaA \nthen \ntrafobA) \nthenhor (\ntrafoaB \nthen \ntrafobB) =
        (\ntrafoaA \nthenhor \ntrafoaB) \nthen (\ntrafobA \nthenhor \ntrafobB)
    \end{equation}
\end{proposition}

\begin{proof}
    Consider the case where we have categories~$\CatC, \CatD, \CatE$, with functors~$\funaA, \funbA\colon\CatC \fto \CatD$ and~$\funaB, \funbB\colon\CatD \fto \CatE$ relating them.
    We can thus have the following four diagrams:
    %
    \equationsag{naturality_interchange_a}{eq:nat-interchange-a}
    %
    \equationsag{naturality_interchange_b}{eq:nat-interchange-b}
    %
    We are now faced with the choice of either vertically composing first, or horizontally composing.
    By vertically composing, we get
    %
    \equationsag{nat-vert-a}{eq:nat-vert-a}
    %
    Subsequently applying horizontal composition yields
    %
    \equationsag{nat-vert-b}{eq:nat-vert-b}
    %
    Otherwise, we apply horizontal composition first:
    %
    \equationsag{nat-hor-vert}{eq:nat-hor-vert}
    %
    ths proving that vertical and horizontal composition are interchangeable.
\end{proof}

\todotext{J: the notation in the proof is inconsistence in itself and with the proposition statement}

\begin{remark}
    The proof is also interesting to do by observing the commuting squares, and putting them together.
\end{remark}



\subsection{Whiskering}
Two special cases of horizontal composition are known as ``left-whiskering'' and ``right-whiskering''. (If you squint your eyes when looking at the diagrams below, and if you, at the same, also think about cats, then probably you will understand the reason for this funny name).


\begin{ctdefinition}[Right Whiskering]\label{def:right-whiskering}
    Let $\CatB, \CatC$ and $\CatD$ be categories, let $\funa, \funb: \CatB \fto \CatC$ and $\func:\CatC \fto \CatD$ be functors, and let $\ntrafoa:\funa \nto \funb$ be a natural transformation. In other words, we have the following situation: 
\equationsag{right_whiskering}{eq:right_whiskering}
The right whiskering of $\ntrafoa$ by $\func$ is the natural transformation
    \begin{equation}
        \func \ntrafoa : \quad (\funa\fthen\func) \quad \nto \quad (\funb\fthen\func)
    \end{equation}
 defined by the following horizontal composition:
 \begin{equation}
        \func \ntrafoa  : = \quad \ntrafoa \nthenhor \natid_\func.
    \end{equation}
    
 In components this means
 \begin{equation}
(\func \ntrafoa)_{\styleobj{B}} =  ({\natid_\func})_{\funa(\styleobj{B})} \mthen \func({\ntrafoa}_{\styleobj{B}}) = \catid_{\func(\funa(\styleobj{B}))} \mthen \func({\ntrafoa}_{\styleobj{B}}) = \func({\ntrafoa}_{\styleobj{B}})
\end{equation}
for any object $\styleobj{B}$ of $\CatB$. 
\end{ctdefinition}



\begin{ctdefinition}[Left Whiskering]\label{def:left-whiskering}
    Let $\CatA, \CatB$ and $\CatC$ be categories, let $\stylefunctors{E}: \CatA \fto \CatB$ and $\funa, \funb: \CatB \fto \CatC$ be functors, and let $\ntrafoa:\funa \nto \funb$ be a natural transformation from $\funa$ to $\funb$. In other words, we are now considering this situation:
    
    \equationsag{left_whiskering}{eq:left_whiskering}

    The left whiskering of $\ntrafoa$ by $\stylefunctors{E}$ is the natural transformation
    \begin{equation}
        \ntrafoa \stylefunctors{E}: (\stylefunctors{E} \fthen\funa) \quad \nto \quad (\stylefunctors{E} \fthen\funb)
    \end{equation}
    defined by the following horizontal composition:
     \begin{equation}
        \ntrafoa \stylefunctors{E}  : = \quad \natid_{\stylefunctors{E}} \nthenhor \ntrafoa.
    \end{equation}
    In components this means
 \begin{equation}
(\stylefunctors{E} \ntrafoa)_{\styleobj{A}} =  {\ntrafoa}_{\stylefunctors{E}(\styleobj{A})} \mthen \funb(({\natid_{\stylefunctors{E}})}_{\styleobj{A}}) =  {\ntrafoa}_{\stylefunctors{E}(\styleobj{A})} \mthen \funb(\catid_{\stylefunctors{E} (\styleobj{A})}) =  {\ntrafoa}_{\stylefunctors{E}(\styleobj{A})} \mthen \catid_{\funb(\stylefunctors{E} (\styleobj{A}))} =  {\ntrafoa}_{\stylefunctors{E}(\styleobj{A})} \end{equation}
for any object $\styleobj{A}$ of $\CatA$. 
\end{ctdefinition}


